\section{The Integer Ladder}

A natural worry is that a chosen honeycomb with a chosen rule is arbitrary. The answer is that the whole substrate
rests on the least arbitrary thing there is, the integers. The Coxeter groups, the crystals, and the cell
addressing all appear as different views of one object built from $\mathbb{Z}$. There are no real numbers and no
tunable dials anywhere in the foundation. The point is best made as a tower, each rung resting on the one below.

\subsection{The tower}

\begin{itemize}
\item \emph{Rung 0, the integers.} Start with $\mathbb{Z}$. Nothing is more canonical. No parameters, no reals,
no randomness.
\item \emph{Rung 1, integer symmetries.} Ask for symmetries generated by reflections where the only data is the
integers saying how the mirrors meet, mirror $i$ and mirror $j$ at angle $\pi$ over an integer. The modular group
$\mathrm{PSL}(2, \mathbb{Z})$, the symmetries of the two-dimensional integer lattice, is one famous example.
\item \emph{Rung 2, Coxeter groups.} That is exactly a \emph{Coxeter group}, and its data is a small integer
matrix, the Coxeter diagram. Still just integers.
\item \emph{Rung 3, tessellations.} Let a Coxeter group act on space by its reflections. Reflect one fundamental
chamber over and over, and the orbit fills the space. The result is a \emph{tessellation}. The algebraic group
and the geometric tiling are the same thing, two views of one object.
\item \emph{Rung 4, the specific crystals.} Choose the integers in the diagram and you get a specific tiling, named
by its Schl\"afli symbol, $\{7,3\}$, $\{5,3,4\}$, $\{3,4,3,4\}$. These are not separate inventions. They are
Rung 3 with particular integer settings, one machine on different dials.
\item \emph{Rung 5, the substrate.} On the chosen tessellation the vibes live, tones on the cells and the rule on
the directions. The substrate inherits, for free, everything the lower rungs provide.
\item \emph{Rung 6, the generating automaton.} Every hyperbolic Coxeter group is an \emph{automatic group}, so a
finite integer state machine generates the whole infinite tessellation deterministically, with no randomness. For
$\{3,4,3,4\}$ this is the solved neighbour automaton and an $O(\log n)$ addressing, by which the substrate
generates and navigates itself from $\mathbb{Z}$ by a finite rule.
\end{itemize}

\subsection{One object, many views}

The power of the ladder is that a single object wears many faces. For $\{3,4,3,4\}$ the Coxeter system
$[3,4,3,4]$ is at once an algebraic group of five reflections, a geometric honeycomb of $24$-cells filling
$\mathbb{H}^4$, a combinatorial addressing tree with Zeckendorf-style digits, a dynamical generator in the
neighbour automaton that advances beat by beat, and a growth law in the shell recurrence. The algebraic Perron
root of that recurrence, near $18.3$, is the central constant of the substrate, the $\{3,4,3,4\}$ analogue of the
golden ratio that sits at the heart of the $\{5,3,4\}$ and heptagrid cases. These are not analogies. They are
literally the same object looked at from several sides.

\subsection{The payoff}

To pick the base is only to pick which integers go in the diagram. The object beneath all the views is one
integer machine, and there is no free real-number parameter anywhere in the foundation.

\begin{proposition}[No free parameters]
$\{3,4,3,4\}$ is not an arbitrary choice. It is one integer setting of the canonical Coxeter machine, and which
setting it is is forced by the selection criteria developed later, the dimension ladder, the crystallographic
spinor coin, and curvature.
\end{proposition}

The integers are the bones of the theory, and the small ones carry the laws. In the song image they are the
score, the small ones its notes, and in the seed image the diagram is the seed that already folds the whole tree.
This is also where the foundation
makes contact with the broader program of generating spacetime from finite combinatorial data, in the spirit of
the automatic-group description of hyperbolic tilings \cite{cannon1997} and the negatively curved geometry that
underlies them \cite{gromov1987}. The substrate is what one gets by taking that machinery seriously as the
ground floor rather than as a tool applied to something else.
